\documentclass[11pt]{amsart}
\usepackage{amsmath,amssymb,amsthm}
\usepackage[margin=1.1in]{geometry}
\usepackage{newtxtext,newtxmath}
\usepackage{microtype}
\usepackage[colorlinks=true, linkcolor=blue, citecolor=blue, urlcolor=blue]{hyperref}
\usepackage{listings}
\usepackage{upquote}
\lstdefinelanguage{Lean}{
  morekeywords={def,theorem,open,in},
  sensitive=true,
  morecomment=[l]{--},
}
\lstset{language=Lean, basicstyle=\small\fontfamily{lmtt}\selectfont\upshape, keywordstyle=\bfseries,
  columns=fullflexible, keepspaces=true, breaklines=false,
  xleftmargin=1em, aboveskip=6pt, belowskip=6pt,
  literate={ℕ}{{$\mathbb{N}$}}1 {ℝ}{{$\mathbb{R}$}}1
    {→}{{$\rightarrow$}}1 {∀}{{$\forall$}}1 {≤}{{$\le$}}1 {≠}{{$\ne$}}1
    {∧}{{$\wedge$}}1 {¬}{{$\neg$}}1 {⌊}{{$\lfloor$}}1 {⌋}{{$\rfloor$}}1
}

\newtheorem{lemma}{Lemma}[section]
\newtheorem{remark}[lemma]{Remark}
\newtheorem*{theoremA}{Theorem}

\newcommand{\Z}{\mathbb{Z}}
\newcommand{\Sset}{\mathcal{S}}
\newcommand{\vt}{\vartheta}

\begin{document}

\title[Prime-prefix-free integers]{The reciprocal sum of the prime-prefix-free integers\\ converges under the Riemann Hypothesis}

\author{Dan Brown$^*$}
\email{jdb1729@gmail.com}
\thanks{$^*$Author of \emph{The Da Vinci Code}.}

\date{October 4, 2026}

\begin{abstract}
Call a positive integer \emph{prime-prefix-free} if none of its proper
binary prefixes $\lfloor n/2^k \rfloor$, $k \ge 1$, is an odd prime (OEIS
A287117). Assuming the Riemann Hypothesis, we prove that the sum of
$1/n$ over the prime-prefix-free integers converges. More precisely, the
proportion of $[2^m, 2^{m+1})$ they occupy is $O_\delta(m^{-1/\log 2 +
\delta})$ for every $\delta > 0$. The heuristic exponent is $1/\log 2
\approx 1.4427$, and convergence rests on the inequality $1/\log 2 > 1$,
that is, on the base $2$ being smaller than $e$. The proof combines
Selberg's theorem on primes in almost all short intervals with an
upper-bound sieve for prime pairs, and closes with a bootstrap for a
delayed recurrence. The convergence statement is machine-verified in
Lean~4 over Mathlib from Mathlib's statement of the Riemann Hypothesis.
Numerically the sum exceeds $3.5276$ unconditionally, and a heuristic
model fitted to exact counts up to $2^{46}$ puts it near $3.877$.
\end{abstract}

\maketitle

\section{Introduction}

Arrange the positive integers in the infinite binary tree in which the
children of $n$ are $2n$ and $2n+1$. The ancestors of $n$ are then its
proper binary prefixes $\lfloor n/2^k \rfloor$, $k \ge 1$. Let $\Sset$ be
the set of $n$ none of whose ancestors is an odd prime. Equivalently, the
odd primes are made leaves of the tree, and $\Sset$ is the set of nodes
that remain. The prime $2$ is excluded because every $n \ge 4$ has
$2$ or $3$ as a prefix, so including it would leave a finite set. The
sequence $\Sset = \{1, 2, 3, 4, 5, 8, 9, 16, 17, 18, 19, 32, \dots\}$ is
A287117 in the OEIS \cite{OEIS}, and the question whether
$\sum_{n \in \Sset} 1/n$ converges was posed on Mathematics Stack
Exchange \cite{MSE}.

\begin{theoremA}
Assume the Riemann Hypothesis. Then $\sum_{n \in \Sset} 1/n < \infty$.
Moreover, for every $\delta > 0$,
\[
\#\bigl(\Sset \cap [2^m, 2^{m+1})\bigr) \;\ll_\delta\; 2^m\, m^{-1/\log 2 + \delta}.
\]
\end{theoremA}

The exponent is the one predicted by a probabilistic model. Pick $x \in
[1,2)$ at random. Its binary prefix of length $m+1$ is an integer near
$2^m$, prime with probability about $1/(m \log 2)$, and these events are
nearly independent, so the prefixes of $x$ avoid the odd primes up to
length $m$ with probability about $\prod_{k \le m}\bigl(1 - 1/(k\log 2)\bigr)
\asymp m^{-1/\log 2}$. That survival probability is the proportion $r_m$
of the level $[2^m, 2^{m+1})$ occupied by $\Sset$, and the reciprocal sum
is comparable to $\sum_m r_m$. Convergence therefore hinges on $1/\log 2
\approx 1.4427$ exceeding $1$. In base $b$ the same model gives the
exponent $1/\log b$, below $1$ for every $b \ge 3$: binary is the one base
in which the analogous sum is expected to converge.

Turning the model into a proof requires lower bounds for primes in the
set of survivors at every level. That set is a union of short intervals
at all scales, so this is a problem about primes in short intervals,
counted simultaneously along all branches of the tree. The proof uses
two known results. The first is Selberg's theorem \cite{Sel}, conditional on RH, that
almost all intervals $[x, x+h]$ with $h$ a small power of $x$ contain the
expected number of primes; the exceptional set is small enough to survive
a union over the polynomially many scales involved. The second is the
classical upper-bound sieve for prime pairs \cite{HR}. It controls primes
that have a prime ancestor only a few levels up. These pairs are beyond
the reach of RH, and the sieve overcounts them by a constant factor, but
they occupy only a small fraction of the levels, and that fraction can be
made as small as we like. A telescoping identity converts the remaining
long-range pairs into the same quantities $r_m$ that the argument
controls, and a bootstrap for the resulting delayed recurrence finishes
the proof.

Each loss in the argument lowers the exponent $1/\log 2$, and the proof
needs the exponent to stay above $1$, so the losses must total less than
$1/\log 2 - 1 \approx 0.44$. In particular, the sieve's constant-factor
overcount can be tolerated only on a small proportion of levels.

\subsection*{Notation}
Logarithms are natural. For $m \ge 0$ let $I_m = [2^m, 2^{m+1}) \cap \Z$,
$\Sset_m = \Sset \cap I_m$, $P_m$ the set of odd primes in $\Sset_m$, and
\[
r_m = |\Sset_m| / 2^m \in [0,1].
\]
For integers $M \ge 1$ and $b \ge 0$ let
$B(M, b) = [M 2^b, (M+1) 2^b) \cap \Z$. These are the descendants of $M$
lying $b$ levels below it: $n \in B(M,b)$ if and only if
$\lfloor n / 2^b \rfloor = M$. For a finite set $A$ of integers write
$\pi(A)$ for the number of primes in $A$ and $\vt(A) = \sum_{p \in A}
\log p$. As usual $\vt(x) = \sum_{p \le x}\log p$. Throughout,
$c = 1/\log 2$.

\section{Components}

\begin{lemma}[Selberg \cite{Sel}]
\label{lem:selberg}
Assume the Riemann Hypothesis. For every $\theta > 0$ there is $C$ such
that for all integers $X \ge 2$ and $h$ with $X^\theta \le h \le X^{1/2}$,
\[
\sum_{X \le n < 2X} \bigl(\vt(n+h) - \vt(n) - h\bigr)^2 \;\le\; C\, h X \log^2 X .
\]
\end{lemma}

Selberg proved the integral form
$\int_X^{2X} (\psi(x+h) - \psi(x) - h)^2\,dx \ll h X \log^2 X$ for
$1 \le h \le X$. For integer $h$ the window $\vt(x+h) - \vt(x)$ is
constant on each $[n, n+1)$, so the integral of the $\vt$-version is the
sum above, and the passage from $\psi$ to $\vt$ costs only the prime
powers $p^k$, $k \ge 2$, which contribute $O(hX\log^2 X)$ to the mean
square. Lemma~\ref{lem:selberg} asks
for much less than Selberg's theorem. Only windows of length at least
$X^\theta$ are needed, for a $\theta$ we will take small but fixed.

\begin{lemma}[Upper-bound sieve for prime pairs; see \cite{HR}]
\label{lem:sieve}
There is an absolute constant $C$ such that for all integers $Y \ge 0$,
$H \ge 2$, $a \ge 1$, $t \ge 1$,
\[
\#\{Y \le n < Y + H : n \text{ and } an + t \text{ prime}\}
\;\le\; C \Bigl(\frac{at}{\varphi(at)}\Bigr)^2 \frac{H}{\log^2 H}.
\]
\end{lemma}

The bound is uniform in $a$ and $t$. The sieve only sifts by primes not
dividing $2at$, and the factor $(at/\varphi(at))^2$ absorbs the
corresponding local densities.

\begin{lemma}
\label{lem:totient}
$\sum_{t \le T} (t/\varphi(t))^2 \ll T$.
\end{lemma}

\begin{proof}
Write $(t/\varphi(t))^2 = \sum_{d \mid t} g(d)$ with $g$ multiplicative,
supported on squarefree $d$, and $g(p) = (p/(p-1))^2 - 1 \le 6/p$. Then the
sum is $\sum_d g(d) \lfloor T/d \rfloor \le T \prod_p (1 + 6/p^2)$.
\end{proof}

\section{The tree}

\begin{lemma}
\label{lem:tree}
\begin{enumerate}
\item[(i)] For $e \in \{0, 1\}$, $2n+e \in \Sset$ if and only if $n \in \Sset$
and $n$ is not an odd prime.
\item[(ii)] If $n \in \Sset$ then $\lfloor n/2^d \rfloor \in \Sset$ for
every $d \ge 0$.
\item[(iii)] $|\Sset_{m+1}| = 2\bigl(|\Sset_m| - |P_m|\bigr)$. Hence
$r_{m+1} = r_m - |P_m|/2^m$, and $r$ is nonincreasing.
\item[(iv)] $\sum_{n \in \Sset} 1/n \le \sum_{m \ge 0} r_m$.
\end{enumerate}
\end{lemma}

\begin{proof}
The proper prefixes of $2n+e$ are $n$ and the proper prefixes of $n$,
which gives (i); (ii) holds because the prefixes of $\lfloor n/2^d\rfloor$
are prefixes of $n$; (iii) follows from (i); and (iv) from
$1/n \le 2^{-m}$ on $I_m$.
\end{proof}

By (iii), $r_m - r_{m+1}$ is the number of odd primes in $\Sset_m$
divided by $2^m$, so upper bounds for $r$ come from lower bounds for these
primes.

\begin{lemma}[Counting inequality]
\label{lem:count}
For integers $j \ge 2$ and $1 \le d \le j$,
\[
\sum_{N \in \Sset_{j-d}} \pi\bigl(B(N, d)\bigr)
\;\le\; |P_j| + \sum_{b=1}^{d}\, \sum_{p \in P_{j-b}} \pi\bigl(B(p, b)\bigr).
\]
\end{lemma}

\begin{proof}
The blocks on the left are disjoint subsets of $I_j$, so the left side
counts primes $q \in I_j$ with $\lfloor q/2^d \rfloor \in \Sset$. Such a
$q$ is at least $2^j \ge 4$, hence odd. If $q \in \Sset$, then $q \in
P_j$. Otherwise some $\lfloor q/2^k\rfloor$ with $k \ge 1$ is an odd
prime. For $k > d$ that prefix is a proper prefix of
$\lfloor q/2^d\rfloor \in \Sset$, so $k \le d$. Let $b \le d$ be the
largest such $k$ and $p = \lfloor q/2^b \rfloor$. Every proper prefix of
$p$ is $\lfloor q/2^{b+k'}\rfloor$ with $k' \ge 1$. For $b + k' \le d$ it
is not an odd prime by the maximality of $b$; for $b + k' > d$ it is a
prefix of an element of $\Sset$. So $p \in P_{j-b}$ and $q \in B(p, b)$.
\end{proof}

Lemma~\ref{lem:count} bounds $|P_j|$ from below by the block counts on
the left, less the sum over $b$, which counts the primes with an
odd-prime ancestor at most $d$ levels up. We split that sum at a
threshold $b_0 = \lceil \theta j \rceil$. Long gaps $b \ge b_0$ are
handled by Lemma~\ref{lem:selberg} in Section~\ref{sec:blocks}, and short
gaps $b < b_0$ by Lemma~\ref{lem:sieve} in Section~\ref{sec:short}.

\section{Blocks}
\label{sec:blocks}

Call a block $B(M, b)$ with $M \in I_{j-b}$ \emph{$\varepsilon$-bad} if
$|\vt(B(M,b)) - 2^b| > \varepsilon 2^b$.

\begin{lemma}[Few bad blocks]
\label{lem:bad}
Assume Lemma~\ref{lem:selberg}. Let $\theta > 0$ and $0 < \varepsilon
\le 1$. There are $C$ and $j_1$ such that for all $j \ge j_1$ and all $b$
with $\theta j \le b \le j/2$, the number of $\varepsilon$-bad blocks
$B(M, b)$, $M \in I_{j-b}$, is at most $C j^3\, 2^j / 4^b$.
\end{lemma}

\begin{proof}
Let $X = 2^j$, $h = 2^b$ and $E(n) = \vt(n+h) - \vt(n) - h$. The
hypotheses give $X^\theta \le h \le X^{1/2}$. For $M \in I_{j-b}$ we have
$\vt(B(M,b)) - h = E(Mh - 1)$. Put $L = \lfloor \varepsilon h / (4 \log
3X) \rfloor$. Since $h \ge X^\theta$, for $j \ge j_1$ the quantity
$\varepsilon h/(4\log 3X)$ is at least $2$, so $L \ge \varepsilon h / (8 \log 3X) \ge 1$. For $Mh \le n < Mh + L$, the
difference $E(n) - E(Mh-1)$ is the $\vt$-mass of two intervals of at
most $L$ integers below $3X$, so it is at most $2L\log 3X \le
\varepsilon h/2$ in absolute value. Hence a bad block yields $L$
consecutive integers $n \in [X, 2X)$ with $|E(n)| > \varepsilon h/2$, and
these runs are disjoint for distinct $M$. By Lemma~\ref{lem:selberg},
\[
\#\{\text{bad } M\} \cdot L \cdot \frac{\varepsilon^2 h^2}{4} \;\le\;
\sum_{X \le n < 2X} E(n)^2 \;\le\; C_1 h X \log^2 X,
\]
so the number of bad blocks is at most
$32 C_1 X \log^2 X \log 3X / (\varepsilon^3 h^2) \ll j^3\, 2^j/4^b$.
\end{proof}

\begin{lemma}[Good blocks]
\label{lem:good}
Let $j \ge 1$, $b \le j$ and $M \in I_{j-b}$. If $B(M,b)$ is not
$\varepsilon$-bad, then
\[
\frac{(1-\varepsilon)\, 2^b}{(j+1) \log 2} \;\le\; \pi\bigl(B(M,b)\bigr)
\;\le\; \frac{(1+\varepsilon)\, 2^b}{j \log 2}.
\]
In every case $\pi(B(M,b)) \le 2^b$.
\end{lemma}

\begin{proof}
$B(M,b) \subset I_j$, so every prime $p$ in it satisfies $j \log 2 \le
\log p < (j+1)\log 2$.
\end{proof}

\section{Short gaps}
\label{sec:short}

\begin{lemma}
\label{lem:short}
Assume Lemma~\ref{lem:sieve}. There is a constant $C_5$ such that for all
integers $1 \le b \le d'$ with $b + d' \le j$,
\[
\sum_{p \in P_{j-b}} \pi\bigl(B(p, b)\bigr)
\;\le\; C_5\, |\Sset_{j-b-d'}|\, \frac{2^{b+d'}}{d'^{\,2}}.
\]
\end{lemma}

\begin{proof}
The primes in $B(p, b)$ are $q = 2^b p + t$ with $0 \le t < 2^b$, and $t$
is odd since $q > 2$. By Lemma~\ref{lem:tree}(ii), $N' = \lfloor
p/2^{d'}\rfloor$ lies in $\Sset_{j-b-d'}$, and $p \in B(N', d')$, an
interval of length $H = 2^{d'}$. Therefore the left side is at most
\[
\sum_{N' \in \Sset_{j-b-d'}}\ \sum_{\substack{t < 2^b\\ t \text{ odd}}}
\#\bigl\{n \in B(N', d') : n \text{ and } 2^b n + t \text{ prime}\bigr\}.
\]
For odd $t$, $2^b t/\varphi(2^b t) = 2t/\varphi(t)$, so
Lemma~\ref{lem:sieve} with $a = 2^b$ bounds each count by
\[
4C \Bigl(\frac{t}{\varphi(t)}\Bigr)^2 \frac{2^{d'}}{(d' \log 2)^2},
\]
and Lemma~\ref{lem:totient} bounds the sum over $t$ by a constant times
$2^b$.
\end{proof}

Heuristically the left side is $|P_{j-b}| \cdot 2^b/(j\log 2)$, which in
terms of $r$ is of order $r\,2^j/j^2$, the same order as the right side
with $d' \asymp j$. The bound is therefore off by a constant factor, and
for this reason the argument applies it only for $b < \theta j$, a small
proportion of the levels.

\section{The recurrence}

\begin{lemma}
\label{lem:rec}
Assume the Riemann Hypothesis. Let $C_5$ be as in Lemma~\ref{lem:short}
and let $0 < \varepsilon \le 1$, $0 < \theta \le 1/7$. There are $K$ and
$j_0$ such that for all $j \ge j_0$,
\[
r_{j+1} \;\le\; r_j \Bigl(1 - \frac{c}{j}\Bigr)
+ \Bigl(\frac{\eta}{j} + \frac{K}{j^2}\Bigr) r_{\lfloor j/3 \rfloor}
+ K j^3\, 2^{-\theta j},
\qquad \eta = 2c\varepsilon + 5 C_5 \theta .
\]
\end{lemma}

\begin{proof}
Let $d = \lfloor j/2 \rfloor$ and $b_0 = \lceil \theta j\rceil$. For $j$
large, $1 \le b_0 \le j/6$, so $j - b - d \ge \lfloor j/3 \rfloor$ for
every $b < b_0$. Apply Lemma~\ref{lem:count} with this $d$ and split the
sum over $b$ at $b_0$.

\emph{Main term.} Every $N \in \Sset_{j-d}$ whose block is not bad
contributes at least $(1-\varepsilon) 2^d/((j+1)\log 2)$ by
Lemma~\ref{lem:good}. Lemma~\ref{lem:bad} with $b = d$ removes at most a
proportion $O(j^3 2^{-d})$ of the $2^{j-d}$ blocks. So the left side of
Lemma~\ref{lem:count}, divided by $2^j$, is at least
$(1-\varepsilon)\, r_{j-d} / ((j+1)\log 2) - O(j^3 2^{-d})$.

\emph{Long gaps, $b_0 \le b \le d$.} By Lemma~\ref{lem:good} for good
blocks, the trivial bound for bad ones, and Lemma~\ref{lem:bad},
\[
\sum_{p \in P_{j-b}} \pi\bigl(B(p,b)\bigr) \;\le\; |P_{j-b}|
\frac{(1+\varepsilon)2^b}{j\log 2} + O\Bigl(\frac{j^3 2^j}{2^b}\Bigr).
\]
By Lemma~\ref{lem:tree}(iii), $|P_{j-b}|\, 2^b = 2^j (r_{j-b} -
r_{j-b+1})$, so the main parts telescope:
\[
\frac{1}{2^j}\sum_{b=b_0}^{d} \sum_{p \in P_{j-b}} \pi\bigl(B(p,b)\bigr)
\;\le\; \frac{(1+\varepsilon)\bigl(r_{j-d} - r_{j-b_0+1}\bigr)}{j \log 2}
+ O\bigl(j^3 2^{-\theta j}\bigr).
\]

\emph{Short gaps, $1 \le b < b_0$.} Lemma~\ref{lem:short} with $d' = d$,
together with $r_{j-b-d} \le r_{\lfloor j/3\rfloor}$, bounds their total,
divided by $2^j$, by $C_5\, b_0\, r_{\lfloor j/3 \rfloor}/d^2 \le (5 C_5
\theta/j + 5C_5/j^2)\, r_{\lfloor j/3\rfloor}$ for $j$ large.

\emph{Assembly.} Write $A = r_{j-d}$ and $B = r_{j-b_0+1} \ge r_j$.
Lemma~\ref{lem:count} and the three estimates give
\[
\frac{|P_j|}{2^j} \;\ge\; \frac{(1-\varepsilon)\frac{j}{j+1} A -
(1+\varepsilon)(A - B)}{j \log 2} - \Bigl(\frac{5C_5\theta}{j} +
\frac{5C_5}{j^2}\Bigr) r_{\lfloor j/3\rfloor} - O\bigl(j^3 2^{-\theta
j}\bigr).
\]
The numerator equals $(1+\varepsilon)B - \bigl(2\varepsilon +
\frac{1-\varepsilon}{j+1}\bigr) A \ge r_j - (2\varepsilon + 1/j)\,
r_{\lfloor j/3 \rfloor}$, using $A \le r_{\lfloor j/3\rfloor}$. Inserting
this into $r_{j+1} = r_j - |P_j|/2^j$ gives the claim.
\end{proof}

In the numerator, the long-gap correction $(1+\varepsilon)(A - B)$
nearly cancels the main term $(1-\varepsilon)A$, leaving about
$r_{j-b_0+1} \ge r_j$. Heuristically, the primes at level $j$ below a
surviving node $d$ levels up number about $r_{j-d}\, 2^j/(j\log 2)$, and
those removed by an earlier prime ancestor account for the decrease of $r$
between levels $j-d$ and $j - b_0$. The cancellation holds up to the
factors $1 \pm \varepsilon$, which is why the long gaps need the
asymptotic count of Lemma~\ref{lem:selberg} and cannot be handled by the
sieve.

\begin{lemma}[Bootstrap]
\label{lem:boot}
Let $r_j \ge 0$, $c, \eta \ge 0$, $\alpha \ge 0$ with $\alpha + 4^\alpha
\eta \le c$, and let $e_j, g_j \ge 0$ with $\sum e_j < \infty$ and $\sum
j^\alpha g_j < \infty$. If
\[
r_{j+1} \;\le\; r_j\Bigl(1 - \frac{c}{j}\Bigr) + \Bigl(\frac{\eta}{j} +
e_j\Bigr) r_{\lfloor j/3\rfloor} + g_j
\]
for all large $j$, then $r_j = O(j^{-\alpha})$.
\end{lemma}

\begin{proof}
Let $s_j = j^\alpha r_j$ and $S_j = \max_{1 \le i \le j} s_i$. For $j
\ge \max(12, c)$ we have $1 - c/j \ge 0$ and $\lfloor j/3 \rfloor \ge
j/4$, so $r_{\lfloor j/3\rfloor} \le 4^\alpha S_j j^{-\alpha}$. Using
$(1+1/j)^\alpha \le e^{\alpha/j}$ and $1 + x \le e^x$,
\[
s_{j+1} \;\le\; S_j\, e^{\alpha/j}\Bigl(1 - \frac{c}{j} +
\frac{4^\alpha\eta}{j} + 4^\alpha e_j\Bigr) + (j+1)^\alpha g_j
\;\le\; S_j\, e^{4^\alpha e_j} + (j+1)^\alpha g_j .
\]
Hence $S_{j+1} \le S_j e^{4^\alpha e_j} + (j+1)^\alpha g_j$, and $S_j$
is bounded.
\end{proof}

\begin{proof}[Proof of the Theorem]
Given $\delta > 0$, set $\alpha = c - \delta$ and choose $\varepsilon$ and
$\theta$ so small that $4^\alpha (2c\varepsilon + 5C_5\theta) \le
\delta$. Lemma~\ref{lem:rec} supplies the hypothesis of
Lemma~\ref{lem:boot} with $e_j = K/j^2$ and $g_j = K j^3 2^{-\theta j}$.
So $|\Sset_m| = 2^m r_m \ll 2^m m^{-c+\delta}$. Taking $\delta < c - 1$,
the sum $\sum r_m$ converges, and Lemma~\ref{lem:tree}(iv) finishes the
proof.
\end{proof}

For the convergence statement alone, the formalization fixes $\alpha =
6/5$, $\varepsilon = 1/200$ and $\theta = 1/(1000(C_5+1))$. Then $\eta
\le 0.02$, $4^{6/5} < 6$ and $6/5 + 6 \cdot 0.02 < 1.44 < c$.

\section{Remarks}

\begin{remark}[Where RH enters]
\label{rem:rh}
The Riemann Hypothesis is used once, through Lemma~\ref{lem:selberg},
and only for windows $h \ge X^\theta$ with $\theta$ small. The
smallness of $\theta$ is forced by the sieve constant: the short-gap
term contributes $5C_5\theta$ to $\eta$, and $\eta$ has to fit inside
the margin $c - 1 \approx 0.44$. Unconditionally,
Lemma~\ref{lem:selberg} is known for $\theta > 1/6$ by Huxley's
zero-density estimate \cite{Hux}, and for $\theta > 2/15$ after Guth and
Maynard \cite{GM}. Both are far too large, since an upper-bound sieve
cannot count the pairs $(p, 2p+1)$, $(p, 4p+t)$, \dots{} within a factor
close to $1$ (the parity problem). A proof with no hypothesis would
therefore need either primes in almost all intervals $[x, x + x^\theta]$
for a small explicit $\theta$, or a way around the short gaps. The
density hypothesis for $\zeta$ would serve in place of RH: it gives
almost-all asymptotics for windows of length $X^\theta$ for every
$\theta > 0$, with exceptional sets small enough for a weaker form of
Lemma~\ref{lem:bad} that the argument tolerates.
\end{remark}

\begin{remark}[Other bases]
In base $b$ the same tree has branching $b$, a prefix at level $m$ is
prime with probability about $1/(m \log b)$, and the model gives $r_m
\approx m^{-1/\log b}$. For $b \ge 3$ the exponent is below $1$ and the
analogous reciprocal sum is expected to diverge. The proof above uses
the size of the exponent only through $c > 1$.
\end{remark}

\begin{remark}[Numerics]
\label{rem:numerics}
We enumerated $\Sset \cap [1, 2^{46})$ exactly. It has $278{,}565{,}951{,}473$
elements, and every odd element at a level below $45$ was tested for
primality. With $\sigma_m = \sum_{n \in \Sset,\, n < 2^{m+1}} 1/n$:
\begin{center}
\begin{tabular}{rrrcc}
$m$ & $|\Sset_m|$ & $|P_m|$ & $m^{1/\log 2}\, r_m$ & $\sigma_m$ \\
\hline
$10$ & $46$ & $7$ & $1.2450$ & $3.156643$ \\
$20$ & $13{,}994$ & $1{,}015$ & $1.0054$ & $3.367433$ \\
$30$ & $7{,}601{,}600$ & $374{,}156$ & $0.9573$ & $3.455907$ \\
$40$ & $5{,}015{,}825{,}530$ & $184{,}008{,}384$ & $0.9342$ & $3.508295$ \\
$44$ & $69{,}489{,}573{,}032$ & $2{,}313{,}408{,}766$ & $0.9281$ & $3.524057$
\end{tabular}
\end{center}
Hence, unconditionally, $\sum_{n \in \Sset} 1/n > \sigma_{45} =
3.527667\ldots$. The normalized density $m^{1/\log 2} r_m$ decreases
slowly. The proportion $h_m = |P_m|/|\Sset_m|$ of odd primes exceeds the
model's prediction $h^0_m = |\Sset_m|^{-1}\sum_{n \in \Sset_m,\, n \text{
odd}} 2/\log n \approx 1/(m \log 2 + 0.057)$ by a factor $1 + 0.53\,h^0_m$,
with the coefficient stable to within $0.02$ over $37 \le m \le 44$. The
excess has a simple source. An odd child $2N+1$ is divisible by an odd
prime $q$ only if $N \equiv (q-1)/2 \not\equiv 0 \pmod q$, and the parents
$N$ that survive have lost the primes, which all lie in nonzero classes.
So the odd elements of $\Sset_m$ are slightly less often divisible by
small primes than random odd integers. An excess proportional to $h^0_m$
changes $r_m$ only by a convergent factor and leaves the exponent
$1/\log 2$ intact.

Extrapolating with $r_{m+1} = r_m(1 - h_m)$, $h_m = (1 + 0.535\,
h^0_m)\, h^0_m$ and $\sum_{n \in \Sset_m} 1/n = 0.9455\, r_m$ (the last
ratio varies by less than $2 \cdot 10^{-4}$ for $m \ge 20$) gives
\[
\sum_{n \in \Sset} \frac{1}{n} \;\approx\; 3.877 .
\]
Starting the extrapolation at any cutoff between $26$ and $45$ moves the
value by less than $0.002$; dropping the excess term gives $3.884$. These
figures are heuristic. The constants in the proof are far too large to
bound the tail rigorously, and the tail beyond $2^{46}$ is still about
$0.35$, since it decays only like $m^{1 - 1/\log 2}$.
\end{remark}

\begin{remark}[Formal verification]
\label{rem:lean}
The convergence statement of the Theorem has been machine-verified in
Lean~4 over Mathlib, through the bound $r_m \ll m^{-6/5}$, from Mathlib's definition of the Riemann Hypothesis and with no
other hypothesis beyond Lean's three standard axioms. The statement,
verbatim from \texttt{PPF/Main.lean} of the formalization (the file
contains nothing else besides imports and a docstring), is:

\noindent\begin{minipage}{\linewidth}
\begin{lstlisting}
def PrimePrefixFree (n : ℕ) : Prop :=
  ∀ k : ℕ, 1 ≤ k → ¬ ((n / 2 ^ k).Prime ∧ n / 2 ^ k ≠ 2)

open Classical in
theorem ppf_summable_of_RH (hRH : RiemannHypothesis) :
    Summable (fun n : ℕ => if PrimePrefixFree n then (1 : ℝ) / n else 0) :=
  (PPF.summable_of_inputs (PPF.selbergMeanSquare_of_RH hRH) PPF.pairSieve).congr
    fun _ => if_congr Iff.rfl rfl rfl
\end{lstlisting}
\end{minipage}

Here \texttt{RiemannHypothesis}, \texttt{Prime} and \texttt{Summable}
are Mathlib's, and natural-number division \texttt{n / 2 \^{} k} is the
floor. The term $n = 0$ contributes $1/0 = 0$ by Lean's convention.

The proof has three audited layers, matching the structure of this
paper. \texttt{summable\_of\_inputs} derives the conclusion from
Lemmas~\ref{lem:selberg} and~\ref{lem:sieve}, stated as hypotheses;
\texttt{pairSieve} proves Lemma~\ref{lem:sieve}; and
\texttt{selbergMeanSquare\_of\_RH} proves Lemma~\ref{lem:selberg} from
RH. The last of these is the main offensive, and its proof follows the
classical route with two adaptations to the available infrastructure.
The truncated explicit formula
\[
\psi(y) = y - \sum_{|\gamma| \le T} \frac{y^\rho}{\rho} + O\Bigl(\frac{y
\log^2 (Ty)}{T} + \log^2(Ty)\Bigr)
\]
(cf.~\cite{Dav}, Ch.~17) is obtained by moving the contour to
$\operatorname{Re} s = -1/2$. On
the left edge and the left halves of the horizontal edges, $\zeta'/\zeta$
is controlled by reflecting through the functional equation, which needs
a bound $|\psi_\Gamma(w)| \ll \log |\operatorname{Im} w|$ for the
digamma function, proved from Euler's limit formula. The mean square of
the zero sum is then computed as in \cite{Sel}, but against the
polynomial weight $\bigl((u - X/2)(4X - u)/X^2\bigr)^2$ majorizing
$1_{[X, 3X]}$. Two integrations by parts turn the kernel
$\int u^{-1 + i(\gamma - \gamma')}\,du$ into $O\bigl(1/(1 + (\gamma -
\gamma')^2)\bigr)$, which with the zero count $N(t+1) - N(t) \ll \log t$
gives $h X \log^2 X$ without the extra logarithm a sharp cutoff would
cost. The sieve generalizes the twin-type bound of the Carmichael
formalization \cite{Bro} to arbitrary intervals and linear forms. That
formalization also supplies the Selberg sieve, the Perron kernel,
contour integrals over rectangles, and zero counting, all reused
verbatim. The new part is about 6{,}800 lines, and the reused part
about 14{,}900. A kernel-checked test confirms that
\texttt{PrimePrefixFree} selects exactly the 61 terms of A287117 below
$536$. The repository is
\url{https://github.com/jdb19937/prime-prefix-free}.
\end{remark}

\begin{thebibliography}{9}

\bibitem{Bro} D. Brown,
\emph{A deterministic quasi-polylogarithmic algorithm for constructing
Carmichael numbers}, 2026.
Formalization: \url{https://github.com/jdb19937/carmichael}.

\bibitem{Dav} H. Davenport,
\emph{Multiplicative Number Theory}, 3rd ed., Graduate Texts in
Mathematics \textbf{74}, Springer, New York, 2000.

\bibitem{GM} L. Guth, J. Maynard,
\emph{New large value estimates for Dirichlet polynomials},
Ann. of Math. (2) \textbf{203} (2026), no.~2.
\href{https://arxiv.org/abs/2405.20552}{arXiv:2405.20552}.

\bibitem{HR} H. Halberstam, H.-E. Richert,
\emph{Sieve Methods}, London Mathematical Society Monographs~\textbf{4},
Academic Press, London, 1974.

\bibitem{Hux} M.~N. Huxley,
\emph{On the difference between consecutive primes},
Invent. Math. \textbf{15} (1972), 164--170.

\bibitem{MSE} D. Brumleve,
\emph{Does the sum of reciprocals of all prime-prefix-free numbers
converge?}, Mathematics Stack Exchange, question 2288648 (2017).
\url{https://math.stackexchange.com/questions/2288648}.

\bibitem{OEIS} OEIS Foundation,
\emph{The On-Line Encyclopedia of Integer Sequences}, sequence A287117.
\url{https://oeis.org/A287117}.

\bibitem{Sel} A. Selberg,
\emph{On the normal density of primes in small intervals, and the
difference between consecutive primes},
Arch. Math. Naturvid. \textbf{47} (1943), no.~6, 87--105.

\end{thebibliography}
\end{document}
